\section{Probe-relative failure categories}
\label{app:failure-grid}

Figure~\ref{fig:failure-grid} summarizes the probe-relative labels used in RQ4. The frozen causal probe supplies the reference decision by thresholding its raw logit at zero. The prompted model decision is obtained from either the normalized Yes/No probability or normalized probabilistic Likert score at threshold 0.5; ties are negative. Labels are assigned separately to every item--prompt observation.

\begin{figure*}[t]
\centering
\large
\renewcommand{\arraystretch}{2.2}
\setlength{\tabcolsep}{16pt}
\resizebox{.92\textwidth}{!}{%
\begin{tabular}{c|c|c}
& \multicolumn{2}{c}{\textbf{Model output}} \\
\textbf{Causal-probe decision} & \textbf{Positive} & \textbf{Negative} \\
\hline
\textbf{Positive}
& \shortstack{\textbf{True positive}\\Probe--output agreement}
& \shortstack{\textbf{Suppression}\\Probe positive, output negative} \\
\hline
\textbf{Negative}
& \shortstack{\textbf{Post-hoc attribution}\\Probe negative, output positive}
& \shortstack{\textbf{True negative}\\Probe--output agreement} \\
\hline
\end{tabular}}
\caption{Visual definition of the probe-relative RQ4 categories. \emph{True positive} and \emph{true negative} denote agreement with the frozen causal-probe decision; they do not assert agreement with dataset ground truth.}
\label{fig:failure-grid}
\end{figure*}
